% GENERATED FILE. Edit canonical lessons, then run npm run generate:books.
% canonical-source-hash: fnv1a64:4fc40082bc031d68
% canonical-lessons: KA-C204-prakatisu, KA-C204-jodisu, KA-C204-ale, KA-C204-maja-madu, KA-C204-spardhisu

\chapter{To publish, To arrange, To measure, To have fun, To compete}
\label{ch:ka-to-publish-to-arrange-to-measure-to-have-fun-to-compete}
\hlchaptermodality{204}

\begin{chapteropening}
\textbf{By the end of this chapter:} \emph{I can say to publish, to arrange, to measure, to have fun and to compete.}

Complete the last lesson of chapter 204: I can say to publish, to arrange, to measure, to have fun and to compete.
\end{chapteropening}

\section[\texorpdfstring{\kn{ಪ್ರಕಟಿಸು} (prakaṭisu)}{ಪ್ರಕಟಿಸು (prakaṭisu)}]{\kn{ಪ್ರಕಟಿಸು} (prakaṭisu) — to publish, to announce}

\label{lesson:KA-C204-prakatisu}



\begin{quote}
Before the new one: say the Kannada for to blow, then the Kannada for to melt.
\end{quote}

\subsection*{You'll want to know: \kn{ಪ್ರಕಟಿಸು}}
\textbf{\kn{ಪ್ರಕಟಿಸು}} — \emph{prakaṭisu} — \textquotedblleft{}to publish, to announce\textquotedblright{}.

\begin{grammarlens}[title={What you've built}]
One more word for this chapter.
\end{grammarlens}

\subsection*{Your turn}
\begin{itemize}
\raggedright
  \item \emph{Say it:} \emph{prakaṭisu}
  \item \emph{Say it:} \emph{prakaṭisu}, once more
  \item \emph{Say it:} say \emph{prakaṭisu}
  \item \emph{Recall:} say the Kannada for to blow, then the Kannada for to melt, then say \emph{prakaṭisu} again
\end{itemize}

\begin{tcolorbox}[breakable,colback=teal!4,colframe=teal!35!black,title={Before you move on}]
What does \kn{ಪ್ರಕಟಿಸು} mean? (\textquotedblleft{}To publish, to announce\textquotedblright{}.) Say it once more.
\end{tcolorbox}
\section[\texorpdfstring{\kn{ಜೋಡಿಸು} (jōḍisu)}{ಜೋಡಿಸು (jōḍisu)}]{\kn{ಜೋಡಿಸು} (jōḍisu) — to arrange, to assemble}

\label{lesson:KA-C204-jodisu}



\begin{quote}
Before the new one: say the Kannada for to melt, then the Kannada for to publish.
\end{quote}

\subsection*{You'll want to know: \kn{ಜೋಡಿಸು}}
\textbf{\kn{ಜೋಡಿಸು}} — \emph{jōḍisu} — \textquotedblleft{}to arrange, to assemble\textquotedblright{}.

\begin{grammarlens}[title={What you've built}]
One more word for this chapter.
\end{grammarlens}

\subsection*{Your turn}
\begin{itemize}
\raggedright
  \item \emph{Say it:} \emph{jōḍisu}
  \item \emph{Say it:} \emph{jōḍisu}, once more
  \item \emph{Say it:} say \emph{jōḍisu}
  \item \emph{Recall:} say the Kannada for to melt, then the Kannada for to publish, then say \emph{jōḍisu} again
\end{itemize}

\begin{tcolorbox}[breakable,colback=teal!4,colframe=teal!35!black,title={Before you move on}]
What does \kn{ಜೋಡಿಸು} mean? (\textquotedblleft{}To arrange, to assemble\textquotedblright{}.) Say it once more.
\end{tcolorbox}
\section[\texorpdfstring{\kn{ಅಳೆ} (aḷe)}{ಅಳೆ (aḷe)}]{\kn{ಅಳೆ} (aḷe) — to measure}

\label{lesson:KA-C204-ale}



\begin{quote}
Before the new one: say the Kannada for to publish, then the Kannada for to arrange.
\end{quote}

\subsection*{You'll want to know: \kn{ಅಳೆ}}
\textbf{\kn{ಅಳೆ}} — \emph{aḷe} — \textquotedblleft{}to measure\textquotedblright{}.

\begin{grammarlens}[title={What you've built}]
One more word for this chapter.
\end{grammarlens}

\subsection*{Your turn}
\begin{itemize}
\raggedright
  \item \emph{Say it:} \emph{aḷe}
  \item \emph{Say it:} \emph{aḷe}, once more
  \item \emph{Say it:} say \emph{aḷe}
  \item \emph{Recall:} say the Kannada for to publish, then the Kannada for to arrange, then say \emph{aḷe} again
\end{itemize}

\begin{tcolorbox}[breakable,colback=teal!4,colframe=teal!35!black,title={Before you move on}]
What does \kn{ಅಳೆ} mean? (\textquotedblleft{}To measure\textquotedblright{}.) Say it once more.
\end{tcolorbox}
\section[\texorpdfstring{\kn{ಮಜಾ} \kn{ಮಾಡು} (majā māḍu)}{ಮಜಾ ಮಾಡು (majā māḍu)}]{\kn{ಮಜಾ} \kn{ಮಾಡು} (majā māḍu) — to have fun}

\label{lesson:KA-C204-maja-madu}



\begin{quote}
Before the new one: say the Kannada for to arrange, then the Kannada for to measure.
\end{quote}

\subsection*{You'll want to know: \kn{ಮಜಾ} \kn{ಮಾಡು}}
\textbf{\kn{ಮಜಾ} \kn{ಮಾಡು}} — \emph{majā māḍu} — \textquotedblleft{}to have fun\textquotedblright{}.

\begin{grammarlens}[title={What you've built}]
One more word for this chapter.
\end{grammarlens}

\subsection*{Your turn}
\begin{itemize}
\raggedright
  \item \emph{Say it:} \emph{majā māḍu}
  \item \emph{Say it:} \emph{majā māḍu}, once more
  \item \emph{Say it:} say \emph{majā māḍu}
  \item \emph{Recall:} say the Kannada for to arrange, then the Kannada for to measure, then say \emph{majā māḍu} again
\end{itemize}

\begin{tcolorbox}[breakable,colback=teal!4,colframe=teal!35!black,title={Before you move on}]
What does \kn{ಮಜಾ} \kn{ಮಾಡು} mean? (\textquotedblleft{}To have fun\textquotedblright{}.) Say it once more.
\end{tcolorbox}
\section[\texorpdfstring{\kn{ಸ್ಪರ್ಧಿಸು} (spardhisu)}{ಸ್ಪರ್ಧಿಸು (spardhisu)}]{\kn{ಸ್ಪರ್ಧಿಸು} (spardhisu) — to compete}

\label{lesson:KA-C204-spardhisu}



\begin{quote}
Before the new one: say the Kannada for to measure, then the Kannada for to have fun.
\end{quote}

\subsection*{You'll want to know: \kn{ಸ್ಪರ್ಧಿಸು}}
\textbf{\kn{ಸ್ಪರ್ಧಿಸು}} — \emph{spardhisu} — \textquotedblleft{}to compete\textquotedblright{}.

\begin{grammarlens}[title={What you've built}]
One more word for this chapter.
\end{grammarlens}

\subsection*{Your turn}
\begin{itemize}
\raggedright
  \item \emph{Say it:} \emph{spardhisu}
  \item \emph{Say it:} \emph{spardhisu}, once more
  \item \emph{Say it:} say \emph{spardhisu}
  \item \emph{Recall:} say the Kannada for to measure, then the Kannada for to have fun, then say \emph{spardhisu} again
\end{itemize}

\begin{tcolorbox}[breakable,colback=teal!4,colframe=teal!35!black,title={Before you move on}]
What does \kn{ಸ್ಪರ್ಧಿಸು} mean? (\textquotedblleft{}To compete\textquotedblright{}.) Say it once more.
\end{tcolorbox}
